% 第十二章：积分变换法

\chapter{积分变换法}

\begin{wulibj}[本章导读]
积分变换法是求解数学物理方程的另一重要方法。通过对自变量进行积分变换，可以将偏微分方程降维，转化为常微分方程或代数方程，大大简化求解过程。本章主要介绍傅里叶变换\index{傅里叶变换}法和拉普拉斯变换\index{拉普拉斯变换}法在数理方程中的应用。

积分变换法的核心思想是：通过变换将复杂问题简化，求解后再通过逆变换回到原问题的解。
\end{wulibj}

\begin{figure}[htbp]
\centering
\begin{tikzpicture}[scale=0.85, every node/.style={font=\small}]
    % 流程框
    \node[draw, rounded corners, fill=blue!10, minimum width=3cm, minimum height=1cm] (pde) at (0, 0) {偏微分方程};
    \node[draw, rounded corners, fill=green!10, minimum width=3cm, minimum height=1cm] (ode) at (5, 0) {常微分方程};
    \node[draw, rounded corners, fill=orange!10, minimum width=3cm, minimum height=1cm] (solve) at (10, 0) {代数方程};
    
    % 箭头
    \draw[->, thick] (pde) -- (ode) node[midway, above] {积分变换};
    \draw[->, thick] (ode) -- (solve) node[midway, above] {再变换};
    
    % 返回箭头
    \draw[->, thick] (solve.south) -- ++(0, -1) -| (pde.south) node[pos=0.25, below] {逆变换};
    
    % 标注
    \node[below=0.3cm of pde] {复杂};
    \node[below=0.3cm of ode] {较简单};
    \node[below=0.3cm of solve] {最简单};
\end{tikzpicture}
\caption{积分变换法的基本思想}
\label{fig:integral-transform-concept}
\end{figure}

% ========== 第一节 ==========
\section{傅里叶变换\index{傅里叶变换}法}

\subsection{基本原理}

傅里叶变换\index{傅里叶变换}可以将偏导数转化为代数运算：
\begin{equation}
    \mathcal{F}\left[\frac{\partial u}{\partial x}\right] = i\omega\hat{u}(\omega), \quad \mathcal{F}\left[\frac{\partial^2 u}{\partial x^2}\right] = -\omega^2\hat{u}(\omega)
\end{equation}

这特别适合处理\textbf{无界区域}的问题。

\begin{figure}[htbp]
\centering
\begin{tikzpicture}[scale=0.8, every node/.style={font=\small}]
    % 左边：空间域
    \begin{scope}[shift={(-4,0)}]
        \draw[->, gray] (-2, 0) -- (2, 0) node[right] {$x$};
        \draw[->, gray] (0, -0.5) -- (0, 2.5) node[above] {$u(x,t)$};
        \draw[blue, thick, domain=-1.8:1.8, samples=100] plot (\x, {1.5*exp(-\x*\x)});
        \node[below] at (0, -1) {空间域（实空间）};
    \end{scope}
    
    % 箭头
    \draw[->, thick, red] (-1.5, 1) -- (1.5, 1) node[midway, above] {$\mathcal{F}$};
    \draw[->, thick, red] (1.5, 0.5) -- (-1.5, 0.5) node[midway, below] {$\mathcal{F}^{-1}$};
    
    % 右边：频域
    \begin{scope}[shift={(4,0)}]
        \draw[->, gray] (-2, 0) -- (2, 0) node[right] {$\omega$};
        \draw[->, gray] (0, -0.5) -- (0, 2.5) node[above] {$\hat{u}(\omega,t)$};
        \draw[blue, thick, domain=-1.8:1.8, samples=100] plot (\x, {1.2*exp(-\x*\x)});
        \node[below] at (0, -1) {频域（变换空间）};
    \end{scope}
\end{tikzpicture}
\caption{傅里叶变换：从空间域到频域}
\label{fig:fourier-transform-domains}
\end{figure}

图 \ref{fig:fourier-transform-domains} 展示了傅里叶变换的本质：将函数从空间域（实空间）变换到频域（变换空间）。在频域中，微分运算变成乘法运算，大大简化了方程。

\subsection{一维热传导方程的傅里叶变换\index{傅里叶变换}解}

考虑无限长细杆的热传导问题：
\begin{equation}
    \begin{cases}
        \dfrac{\partial u}{\partial t} = a^2\dfrac{\partial^2 u}{\partial x^2}, & x \in \mathbb{R}, \; t > 0 \\
        u(x, 0) = \varphi(x), & x \in \mathbb{R}
    \end{cases}
\end{equation}

\textbf{求解步骤}：

\textbf{第一步}：对 $x$ 作傅里叶变换\index{傅里叶变换}

设 $\hat{u}(\omega, t) = \mathcal{F}[u(x, t)]$，方程变为：
\begin{equation}
    \frac{\partial \hat{u}}{\partial t} = -a^2\omega^2\hat{u}
\end{equation}

初始条件：
\begin{equation}
    \hat{u}(\omega, 0) = \hat{\varphi}(\omega)
\end{equation}

\textbf{第二步}：求解常微分方程

\begin{equation}
    \hat{u}(\omega, t) = \hat{\varphi}(\omega)\mathrm{e}^{-a^2\omega^2 t}
\end{equation}

\textbf{第三步}：求傅里叶逆变换

利用卷积定理：
\begin{equation}
    u(x, t) = \varphi(x) * \mathcal{F}^{-1}[\mathrm{e}^{-a^2\omega^2 t}]
\end{equation}

需要计算 $\mathcal{F}^{-1}[\mathrm{e}^{-a^2\omega^2 t}]$。

利用高斯函数的傅里叶变换\index{傅里叶变换}：
\begin{equation}
    \mathcal{F}[\mathrm{e}^{-\alpha x^2}] = \sqrt{\frac{\pi}{\alpha}}\mathrm{e}^{-\omega^2/(4\alpha)}
\end{equation}

令 $\alpha = \frac{1}{4a^2t}$，得
\begin{equation}
    \mathcal{F}^{-1}[\mathrm{e}^{-a^2\omega^2 t}] = \frac{1}{2a\sqrt{\pi t}}\mathrm{e}^{-x^2/(4a^2t)}
\end{equation}

\begin{dingli}{泊松积分公式}{}
一维热传导方程初值问题的解为
\[
    u(x, t) = \frac{1}{2a\sqrt{\pi t}}\int_{-\infty}^{+\infty} \varphi(\xi)\mathrm{e}^{-(x-\xi)^2/(4a^2t)}\,\mathrm{d}\xi
\]
\end{dingli}

\begin{figure}[htbp]
\centering
\begin{tikzpicture}[scale=0.75]
    % 泊松核 K(x,t) = exp(-x^2/(4t)) / sqrt(4*pi*t)
    % t=0.1, 0.5, 1.0, 2.0
    \draw[->, gray] (-5, 0) -- (5, 0) node[right] {$x$};
    \draw[->, gray] (0, 0) -- (0, 4) node[above] {$K(x,t)$};
    
    % t=0.1 (highest, narrowest)
    \draw[blue, thick, domain=-2:2, samples=100] plot (\x, {3*exp(-\x*\x/0.4)});
    \node[blue] at (2.5, 2.5) {\small $t=0.1$};
    
    % t=0.5
    \draw[red, thick, domain=-3:3, samples=100] plot (\x, {1.5*exp(-\x*\x/2)});
    \node[red] at (3.2, 1.2) {\small $t=0.5$};
    
    % t=1.0
    \draw[green!60!black, thick, domain=-4:4, samples=100] plot (\x, {1*exp(-\x*\x/4)});
    \node[green!60!black] at (3.8, 0.6) {\small $t=1.0$};
    
    % t=2.0
    \draw[orange, thick, domain=-4.5:4.5, samples=100] plot (\x, {0.7*exp(-\x*\x/8)});
    \node[orange] at (4.2, 0.35) {\small $t=2.0$};
    
    % 标注
    \node at (0, -1) {\small 泊松核 $K(x,t) = \dfrac{1}{2a\sqrt{\pi t}}\mathrm{e}^{-x^2/(4a^2t)}$};
\end{tikzpicture}
\caption{泊松核在不同时刻的形状（扩散过程）}
\label{fig:poisson-kernel}
\end{figure}

从图 \ref{fig:poisson-kernel} 可以看出泊松核的性质：
\begin{itemize}
    \item 当 $t \to 0^+$ 时，泊松核趋于 Dirac $\delta$ 函数（初始条件的影响）
    \item 当 $t$ 增大时，核函数变宽变矮（热量扩散到更大范围）
    \item 核函数的积分恒为 1（热量守恒）
\end{itemize}

\begin{lizi}{}{}
求解：
\begin{equation}
    \begin{cases}
        u_t = u_{xx}, \quad x \in \mathbb{R}, \; t > 0 \\
        u(x, 0) = \mathrm{e}^{-x^2}
    \end{cases}
\end{equation}

\begin{jie}
    由泊松积分公式：
    \begin{equation}
        u(x, t) = \frac{1}{2\sqrt{\pi t}}\int_{-\infty}^{+\infty} \mathrm{e}^{-\xi^2}\mathrm{e}^{-(x-\xi)^2/(4t)}\,\mathrm{d}\xi
    \end{equation}
    
    另一种方法：先作傅里叶变换\index{傅里叶变换}。
    \begin{equation}
        \hat{\varphi}(\omega) = \mathcal{F}[\mathrm{e}^{-x^2}] = \sqrt{\pi}\mathrm{e}^{-\omega^2/4}
    \end{equation}
    
    所以
    \begin{equation}
        \hat{u}(\omega, t) = \sqrt{\pi}\mathrm{e}^{-\omega^2/4}\mathrm{e}^{-\omega^2 t} = \sqrt{\pi}\mathrm{e}^{-\omega^2(1/4+t)}
    \end{equation}
    
    逆变换：
    \begin{equation}
        u(x, t) = \sqrt{\frac{\pi}{1/4+t}}\mathrm{e}^{-x^2/(4t+1)} = \frac{1}{\sqrt{4t+1}}\mathrm{e}^{-x^2/(4t+1)}
    \end{equation}
\end{jie}
\end{lizi}

\subsection{波动方程的傅里叶变换\index{傅里叶变换}解}

对于一维波动方程初值问题：
\begin{equation}
    \begin{cases}
        u_{tt} = a^2u_{xx}, & x \in \mathbb{R}, \; t > 0 \\
        u(x, 0) = \varphi(x), \quad u_t(x, 0) = \psi(x)
    \end{cases}
\end{equation}

对 $x$ 作傅里叶变换\index{傅里叶变换}：
\begin{equation}
    \frac{\partial^2 \hat{u}}{\partial t^2} + a^2\omega^2\hat{u} = 0
\end{equation}

解为
\begin{equation}
    \hat{u}(\omega, t) = \hat{\varphi}(\omega)\cos a\omega t + \frac{\hat{\psi}(\omega)}{a\omega}\sin a\omega t
\end{equation}

利用傅里叶逆变换可回到达朗贝尔公式。

% ========== 第二节 ==========
\section{拉普拉斯变换\index{拉普拉斯变换}法}

\subsection{基本原理}

拉普拉斯变换\index{拉普拉斯变换}特别适合处理\textbf{半无界问题}和\textbf{初值问题}：
\begin{equation}
    \mathcal{L}\left[\frac{\partial u}{\partial t}\right] = s\tilde{u}(x, s) - u(x, 0)
\end{equation}

\subsection{半无界热传导问题}

考虑半无界细杆的热传导问题：
\begin{equation}
    \begin{cases}
        u_t = a^2u_{xx}, & x > 0, \; t > 0 \\
        u(x, 0) = 0 \\
        u(0, t) = f(t), \quad u(+\infty, t) \text{ 有界}
    \end{cases}
\end{equation}

\textbf{求解步骤}：

\textbf{第一步}：对 $t$ 作拉普拉斯变换\index{拉普拉斯变换}

设 $\tilde{u}(x, s) = \mathcal{L}[u(x, t)]$，方程变为：
\begin{equation}
    s\tilde{u} = a^2\frac{\partial^2 \tilde{u}}{\partial x^2}
\end{equation}

边界条件：
\begin{equation}
    \tilde{u}(0, s) = \tilde{f}(s), \quad \tilde{u}(+\infty, s) \text{ 有界}
\end{equation}

\textbf{第二步}：求解常微分方程

\begin{equation}
    \tilde{u}(x, s) = A\mathrm{e}^{\sqrt{s}x/a} + B\mathrm{e}^{-\sqrt{s}x/a}
\end{equation}

由有界性条件，$A = 0$。

由 $\tilde{u}(0, s) = \tilde{f}(s)$ 得 $B = \tilde{f}(s)$。

所以
\begin{equation}
    \tilde{u}(x, s) = \tilde{f}(s)\mathrm{e}^{-\sqrt{s}x/a}
\end{equation}

\textbf{第三步}：求拉普拉斯逆变换

利用卷积定理：
\begin{equation}
    u(x, t) = f(t) * \mathcal{L}^{-1}[\mathrm{e}^{-\sqrt{s}x/a}]
\end{equation}

已知
\begin{equation}
    \mathcal{L}^{-1}[\mathrm{e}^{-\sqrt{s}x/a}] = \frac{x}{2a\sqrt{\pi t^3}}\mathrm{e}^{-x^2/(4a^2t)}
\end{equation}

所以
\begin{equation}
    \boxed{u(x, t) = \int_0^t f(\tau)\frac{x}{2a\sqrt{\pi(t-\tau)^3}}\mathrm{e}^{-x^2/[4a^2(t-\tau)]}\,\mathrm{d}\tau}
\end{equation}

\begin{lizi}{}{}
求解：
\begin{equation}
    \begin{cases}
        u_t = u_{xx}, & x > 0, \; t > 0 \\
        u(x, 0) = 0 \\
        u(0, t) = T_0 \quad \text{（常数）}
    \end{cases}
\end{equation}

\begin{jie}
    $\tilde{f}(s) = T_0/s$。
    
    \begin{equation}
        \tilde{u}(x, s) = \frac{T_0}{s}\mathrm{e}^{-\sqrt{s}x}
    \end{equation}
    
    利用积分性质 $\mathcal{L}^{-1}[\frac{1}{s}F(s)] = \int_0^t f(\tau)\,\mathrm{d}\tau$：
    \begin{equation}
        u(x, t) = T_0\int_0^t \frac{x}{2\sqrt{\pi\tau^3}}\mathrm{e}^{-x^2/(4\tau)}\,\mathrm{d}\tau
    \end{equation}
    
    令 $\eta = x/(2\sqrt{\tau})$，可化为
    \begin{equation}
        u(x, t) = T_0\operatorname{erfc}\left(\frac{x}{2\sqrt{t}}\right) = T_0\left[1 - \operatorname{erf}\left(\frac{x}{2\sqrt{t}}\right)\right]
    \end{equation}
    
    其中 $\operatorname{erf}(x) = \frac{2}{\sqrt{\pi}}\int_0^x \mathrm{e}^{-\eta^2}\,\mathrm{d}\eta$ 是误差函数\index{误差函数}。
\end{jie}
\end{lizi}

\begin{figure}[htbp]
\centering
\begin{tikzpicture}[scale=0.8]
    % 坐标轴
    \draw[->, gray] (-0.5, 0) -- (5, 0) node[right] {$x$};
    \draw[->, gray] (0, -0.2) -- (0, 3.5) node[above] {$u(x,t)/T_0$};
    
    % 温度分布曲线（示意图）
    % t=0.01: 快速下降
    \draw[blue, thick, smooth] plot coordinates {(0,3) (0.3,2.8) (0.6,2) (1,0.8) (1.5,0.2) (2,0.05) (2.5,0)};
    \node[blue] at (1.5, 2.5) {\small $t=0.01$};
    
    % t=0.1: 较慢下降
    \draw[red, thick, smooth] plot coordinates {(0,3) (0.5,2.7) (1,2) (1.5,1.2) (2,0.6) (2.5,0.3) (3,0.1) (3.5,0)};
    \node[red] at (2.8, 1.8) {\small $t=0.1$};
    
    % t=0.5: 更慢下降
    \draw[green!60!black, thick, smooth] plot coordinates {(0,3) (0.7,2.6) (1.5,1.8) (2,1.2) (2.5,0.8) (3,0.5) (3.5,0.3) (4,0.15) (4.5,0)};
    \node[green!60!black] at (3.5, 1.2) {\small $t=0.5$};
    
    % t=1.0: 最慢下降
    \draw[orange, thick, smooth] plot coordinates {(0,3) (1,2.4) (1.5,1.8) (2,1.3) (2.5,0.9) (3,0.65) (3.5,0.45) (4,0.3) (4.5,0.2)};
    \node[orange] at (4, 0.7) {\small $t=1.0$};
    
    % 边界条件标注
    \draw[dashed] (0, 1) -- (5, 1);
    \node[left] at (0, 1) {$1$};
    \node[left] at (0, 0) {$0$};
    
    % 说明
    \node at (2.5, -0.8) {\small 半无界热传导问题的解：$u(x,t) = T_0 \cdot \mathrm{erfc}\!\left(\dfrac{x}{2\sqrt{t}}\right)$};
\end{tikzpicture}
\caption{半无界热传导问题的温度分布随时间的演化}
\label{fig:semi-infinite-heat}
\end{figure}

图 \ref{fig:semi-infinite-heat} 展示了半无界细杆的温度分布随时间的演化：
\begin{itemize}
    \item 边界 $x=0$ 处保持恒温 $T_0$
    \item 热量逐渐向 $x>0$ 区域扩散
    \item 时间越长，温度前沿传播越远
\end{itemize}

\subsection{拉普拉斯变换\index{拉普拉斯变换}解常微分方程}

拉普拉斯变换\index{拉普拉斯变换}也可用于求解常微分方程初值问题，这在分离变量后处理时间方程时很有用。

\begin{lizi}{}{}
求解：
\begin{equation}
    y'' + \omega^2 y = f(t), \quad y(0) = y'(0) = 0
\end{equation}

\begin{jie}
    作拉普拉斯变换\index{拉普拉斯变换}：
    \begin{equation}
        (s^2 + \omega^2)\tilde{y}(s) = \tilde{f}(s)
    \end{equation}
    
    所以
    \begin{equation}
        \tilde{y}(s) = \frac{\tilde{f}(s)}{s^2 + \omega^2}
    \end{equation}
    
    逆变换：
    \begin{equation}
        y(t) = \frac{1}{\omega}\int_0^t f(\tau)\sin\omega(t-\tau)\,\mathrm{d}\tau
    \end{equation}
\end{jie}
\end{lizi}

% ========== 第三节 ==========
\section{积分变换法与分离变量法\index{分离变量法}的比较}

\begin{table}[htbp]
\centering
\begin{tabular}{lll}
\toprule
\textbf{方法} & \textbf{适用问题} & \textbf{特点} \\
\midrule
傅里叶变换\index{傅里叶变换} & 无界区域问题 & 自动满足无穷远条件 \\
拉普拉斯变换\index{拉普拉斯变换} & 半无界、初值问题 & 自动处理初始条件 \\
分离变量法\index{分离变量法} & 有界区域问题 & 得到无穷级数解 \\
\bottomrule
\end{tabular}
\caption{三种方法的比较}
\end{table}

\begin{jinshi}[方法选择]
\begin{itemize}
    \item 无界区域 $\Rightarrow$ 傅里叶变换\index{傅里叶变换}
    \item 半无界区域 $\Rightarrow$ 拉普拉斯变换\index{拉普拉斯变换}（对时间变量）
    \item 有界区域 $\Rightarrow$ 分离变量法\index{分离变量法}
    \item 非齐次方程 $\Rightarrow$ 积分变换法（更方便）
\end{itemize}
\end{jinshi}

% ========== 本章小结 ==========
\section{本章小结}

\subsection{傅里叶变换\index{傅里叶变换}法}

适用：无界区域问题

核心步骤：
\begin{enumerate}
    \item 对空间变量作傅里叶变换\index{傅里叶变换}
    \item 求解关于变换变量的常微分方程
    \item 作傅里叶逆变换
\end{enumerate}

结果：泊松积分公式（热传导）、达朗贝尔公式（波动）

\subsection{拉普拉斯变换\index{拉普拉斯变换}法}

适用：半无界问题、初值问题

核心步骤：
\begin{enumerate}
    \item 对时间变量作拉普拉斯变换\index{拉普拉斯变换}
    \item 求解关于空间变量的常微分方程
    \item 作拉普拉斯逆变换
\end{enumerate}

常用技巧：
\begin{itemize}
    \item 卷积定理
    \item 积分性质
    \item 查表法
\end{itemize}

% ========== 习题 ==========
\section{习题}

\textbf{基础题}

\begin{enumerate}
    \item 用傅里叶变换\index{傅里叶变换}法求解：
    \begin{equation}
        \begin{cases}
            u_t = u_{xx}, \quad x \in \mathbb{R}, \; t > 0 \\
            u(x, 0) = \begin{cases}
                1, & |x| < 1 \\
                0, & |x| > 1
            \end{cases}
        \end{cases}
    \end{equation}
    
    \item 用拉普拉斯变换\index{拉普拉斯变换}法求解：
    \begin{equation}
        \begin{cases}
            u_t = u_{xx}, \quad x > 0, \; t > 0 \\
            u(x, 0) = 0 \\
            u(0, t) = \sin\omega t
        \end{cases}
    \end{equation}
    
    \item 用傅里叶变换\index{傅里叶变换}法求解热传导方程，设初值为高斯型分布：
    \begin{equation}
        \begin{cases}
            u_t = a^2 u_{xx}, \quad x \in \mathbb{R}, \; t > 0 \\
            u(x, 0) = A\mathrm{e}^{-x^2/\sigma^2}
        \end{cases}
    \end{equation}
    其中 $A$ 和 $\sigma$ 为正常数。讨论解随时间的演化特征。
    
    \item 用拉普拉斯变换\index{拉普拉斯变换}法求解半无界热传导问题：
    \begin{equation}
        \begin{cases}
            u_t = a^2 u_{xx}, \quad x > 0, \; t > 0 \\
            u(x, 0) = u_0 \\
            u(0, t) = 0
        \end{cases}
    \end{equation}
    并解释解的物理意义。
    
    \item 用傅里叶变换\index{傅里叶变换}法求解：
    \begin{equation}
        \begin{cases}
            u_{tt} = a^2 u_{xx}, \quad x \in \mathbb{R}, \; t > 0 \\
            u(x, 0) = 0, \quad u_t(x, 0) = \delta(x)
        \end{cases}
    \end{equation}
    验证结果与达朗贝尔公式\index{达朗贝尔公式}一致。
    
    \item 用拉普拉斯变换\index{拉普拉斯变换}法求解半无界杆的热传导问题：
    \begin{equation}
        \begin{cases}
            u_t = a^2 u_{xx}, \quad x > 0, \; t > 0 \\
            u(x, 0) = 0 \\
            u_x(0, t) = g(t), \quad \lim_{x\to+\infty} u(x, t) < \infty
        \end{cases}
    \end{equation}
\end{enumerate}

\textbf{综合题}

\begin{enumerate}[resume]
    \item 用傅里叶变换\index{傅里叶变换}法求解一维波动方程初值问题，验证结果与达朗贝尔公式一致。
    
    \item 求解半无界弦的振动问题：
    \begin{equation}
        \begin{cases}
            u_{tt} = a^2u_{xx}, \quad x > 0, \; t > 0 \\
            u(x, 0) = 0, \quad u_t(x, 0) = 0 \\
            u(0, t) = f(t)
        \end{cases}
    \end{equation}
    
    \item 用傅里叶变换\index{傅里叶变换}法求解一维波动方程初值问题：
    \begin{equation}
        \begin{cases}
            u_{tt} = a^2 u_{xx}, \quad x \in \mathbb{R}, \; t > 0 \\
            u(x, 0) = \varphi(x), \quad u_t(x, 0) = \psi(x)
        \end{cases}
    \end{equation}
    其中 $\varphi(x)$ 和 $\psi(x)$ 有紧支集。
    
    \item 用拉普拉斯变换\index{拉普拉斯变换}法求解非齐次热传导方程的半无界问题：
    \begin{equation}
        \begin{cases}
            u_t = a^2 u_{xx} + f(x,t), \quad x > 0, \; t > 0 \\
            u(x, 0) = 0 \\
            u(0, t) = 0, \quad \lim_{x\to+\infty} u(x, t) < \infty
        \end{cases}
    \end{equation}
    
    \item 求解非齐次波动方程的初值问题：
    \begin{equation}
        \begin{cases}
            u_{tt} = a^2 u_{xx} + h(x,t), \quad x \in \mathbb{R}, \; t > 0 \\
            u(x, 0) = 0, \quad u_t(x, 0) = 0
        \end{cases}
    \end{equation}
    试用傅里叶变换\index{傅里叶变换}法求解，并与推迟势\index{推迟势}公式比较。
    
    \item 用积分变换法求解含源热传导方程：
    \begin{equation}
        \begin{cases}
            u_t = a^2 u_{xx} + Q_0\delta(x), \quad x \in \mathbb{R}, \; t > 0 \\
            u(x, 0) = 0
        \end{cases}
    \end{equation}
    解释解的物理意义。
\end{enumerate}

\textbf{挑战题}

\begin{enumerate}[resume]
    \item 用积分变换法求解二维拉普拉斯方程在上半平面的狄利克雷问题：
    \begin{equation}
        \begin{cases}
            u_{xx} + u_{yy} = 0, \quad y > 0 \\
            u(x, 0) = f(x)
        \end{cases}
    \end{equation}
    结果应为泊松积分公式：
    \begin{equation}
        u(x, y) = \frac{y}{\pi}\int_{-\infty}^{+\infty} \frac{f(\xi)}{(x-\xi)^2 + y^2}\,\mathrm{d}\xi
    \end{equation}
    
    \item 用傅里叶变换\index{傅里叶变换}法求解三维波动方程的初值问题：
    \begin{equation}
        \begin{cases}
            u_{tt} = a^2 \Delta u, \quad \mathbf{r} \in \mathbb{R}^3, \; t > 0 \\
            u(\mathbf{r}, 0) = \varphi(\mathbf{r}), \quad u_t(\mathbf{r}, 0) = \psi(\mathbf{r})
        \end{cases}
    \end{equation}
    导出三维波动方程的泊松公式\index{泊松公式}。
    
    \item 用拉普拉斯变换\index{拉普拉斯变换}法求解具有第三类边界条件的热传导问题：
    \begin{equation}
        \begin{cases}
            u_t = a^2 u_{xx}, \quad x > 0, \; t > 0 \\
            u(x, 0) = f(x) \\
            u_x(0, t) - hu(0, t) = 0, \quad h > 0
        \end{cases}
    \end{equation}
    讨论边界条件参数 $h$ 对解的影响。
    
    \item 考虑对流扩散方程的初值问题：
    \begin{equation}
        \begin{cases}
            u_t + v u_x = D u_{xx}, \quad x \in \mathbb{R}, \; t > 0 \\
            u(x, 0) = \varphi(x)
        \end{cases}
    \end{equation}
    其中 $v > 0$ 为对流速度，$D > 0$ 为扩散系数。用傅里叶变换\index{傅里叶变换}法求解，分析对流项与扩散项对解的影响。
\end{enumerate}
